\documentclass[11pt]{article}
\usepackage[margin=1in]{geometry}
\usepackage{amsmath}
\usepackage{booktabs}
\usepackage{graphicx}
\usepackage{hyperref}

\title{Adjoint Verification for the Dissipative-Wave Inverse Problem\\
Case 3 (High Contrast) Reference Study}
\author{REMAT Example Study}
\date{April 28, 2026}

\begin{document}
\maketitle

\begin{abstract}
This report verifies adjoint-gradient correctness for the dissipative-wave inverse problem, using \texttt{case3\_high\_contrast} as the reference configuration.
Four levels of checks were performed: (i) existing C++ unit tests for material- and system-level adjoint kernels, (ii) a one-step material-point transpose-identity test for constitutive reverse-mode consistency, (iii) an element-level coupling audit on a small fixed-point mesh, and (iv) finite-difference (FD) checks in the full high-contrast inverse setup.
Unit tests pass; however, in the full Case 3 setup (fixed-point visco integrator with frequent overflow checkpointing), layer-parameter gradients show persistent FD mismatch, while $\tau$ agreement is step-size dependent.
Windowed directional checks further show that directional inconsistency is injected predominantly in mid/late parts of the trajectory rather than uniformly over time.
An additional reduced-stiffness/reduced-horizon control test with a float-visco auxiliary cross-check shows strong fixed-vs-float gradient-direction agreement and moderate FD consistency.
Therefore, we cannot claim full adjoint correctness for Case 3 layer gradients under the current numerical regime, but the control-regime evidence supports the adjoint implementation.
\end{abstract}

\section{Case 3 Reference Setup}
We used the preset in \texttt{examples/dissipative\_wave/dissipative\_wave\_inverse\_cases.py}:
\begin{itemize}
\item Domain: width \(=15\), height \(=3\)
\item Mesh: \(150\times 30\) quads
\item Time integration: \(N_t=1000\), \(\Delta t=4\times 10^{-3}\)
\item Parameters: \(s=[s_0,s_1]\), \(\tau\), with true \(s=[15,7.7]\), true \(\tau=0.08\)
\item Initial guess: \(s=[10,4]\), \(\tau=0.10\)
\item Sensors: 5 top-surface velocity-\(y\) sensors
\item Integrator: \texttt{fixed\_visco}
\item Overflow settings: \texttt{overflow\_limit=10}, \texttt{mat\_overflow\_limit=10}
\end{itemize}

The objective is
\[
J(\theta)=\frac12\sum_{k=0}^{N_t-1}\left\|v_y^{\mathrm{sim}}(t_k;\theta)-v_y^{\mathrm{obs}}(t_k)\right\|_2^2,
\quad
\theta=[s_0,s_1,\tau]^T.
\]
Component FD checks use central differences:
\[
\frac{\partial J}{\partial \theta_i}\approx
\frac{J(\theta+h e_i)-J(\theta-h e_i)}{2h}.
\]
Directional checks use
\[
\nabla J(\theta)\cdot d \approx \frac{J(\theta+h d)-J(\theta-h d)}{2h}.
\]

\subsection*{Why Fixed-Point is Used, and Why Tiny FD Steps are Avoided}
The use of fixed-point arithmetic is intentional: it enables reverse-time rematerialization of integration history instead of relying on heavy checkpointing schemes.
In this project, increasing overflow limits tends to accumulate roundoff; with damping present, that accumulated roundoff can wipe out informative signal in forward/adjoint histories and destabilize the inversion dynamics.
Float-visco can still be useful as an auxiliary check in reduced-horizon, larger-\(\tau\), lower-stiffness controls, but it is not the primary production mode for the high-contrast rematerialized runs.
Also, in \texttt{fixed\_visco}, sufficiently small perturbations can become numerically unresolved by the fixed precision and effectively vanish.
Accordingly, we use \(h=10^{-2}\) for finite-difference verification as a practical balance: small enough to probe local sensitivity, yet large enough to remain above fixed-point resolution and avoid roundoff-dominated cancellation.

\section{Baseline Kernel Verification}
We first ran targeted C++ tests:
\begin{itemize}
\item \texttt{test\_UniaxialViscoelasticity}
\item \texttt{test\_ViscoElasticityAdjoint}
\item \texttt{test\_SystemAdjointKick}
\end{itemize}
All passed:
\begin{center}
\texttt{100\% tests passed, 0 tests failed out of 3}
\end{center}
This indicates that core adjoint operators are internally consistent on their unit-test configurations.

\section{Material-Point Transpose Validation}
To isolate constitutive reverse-mode consistency from global assembly effects, we performed a one-step material-point transpose test for the viscoelastic history update.
At fixed strain state and time step, we considered the local map
\[
q_n \mapsto y(q_n)=\big[q_{n+1}(q_n),\ \sigma_n(q_n)\big],
\]
and evaluated the discrete adjoint identity
\[
w^\top\!\left(J\,v\right)=\left(J^\top w\right)^\top v,
\qquad J=\frac{\partial y}{\partial q_n},
\]
for independent perturbation directions \(v\) and cotangent directions \(w\).

The Jacobian-vector product \(Jv\) was estimated by centered finite differences of the local forward/rematerialized map, while \(J^\top w\) was obtained from the implemented material pullback in reverse mode.
Across four independent direction-pair checks, the maximum absolute mismatch was \(7.22\times10^{-13}\), the maximum relative mismatch was \(2.38\times10^{-12}\), and the mean relative mismatch was \(1.41\times10^{-12}\).
These machine-precision closures indicate that the constitutive adjoint and the local coupling between stress seeds and history-state seeds satisfy the expected discrete transpose relation in the tested regime.

\section{Element-Level Coupling Audit}
To verify coupling between constitutive sensitivities and the global objective accumulation, we performed a small-mesh audit with element-wise perturbations.
The test used a \(6\times2\) quadrilateral mesh, fixed-point visco integration, \(\tau=1\), \(N_t=220\), and \(\Delta t=2.5\times10^{-3}\).
Element stiffness-scaling factors were treated as independent local controls for the audit, and the objective residual was formed from top-surface velocity sensors.

Let \(g^{\mathrm{adj}}_e\) denote the adjoint sensitivity field \(\partial J/\partial s_e\) returned at each element \(e\), and let \(g^{\mathrm{FD}}_e\) denote centered finite-difference sensitivities obtained by perturbing only element \(e\) with \(h=10^{-2}\).
We evaluated agreement using cosine similarity and a relative \(L^2\) mismatch,
\[
\varepsilon_{L^2}=\frac{\|g^{\mathrm{adj}}-g^{\mathrm{FD}}\|_2}
{\max(\|g^{\mathrm{adj}}\|_2,\|g^{\mathrm{FD}}\|_2,10^{-14})}.
\]

The full-field comparison gave
\[
\cos(g^{\mathrm{adj}},g^{\mathrm{FD}})=0.9999993,\qquad
\varepsilon_{L^2}=7.93\times10^{-3}.
\]
Because many elements are weakly active in this control, component-wise relative errors are inflated by near-zero denominators.
Accordingly, we also report active-element statistics: for elements with \(\max(|g^{\mathrm{adj}}_e|,|g^{\mathrm{FD}}_e|)\ge 10^{-3}\), the relative \(L^2\) mismatch is \(7.92\times10^{-3}\) with \(100\%\) sign agreement; for threshold \(10^{-2}\), the mismatch is \(7.84\times10^{-3}\), also with \(100\%\) sign agreement.

These results indicate that element-level sensitivity transfer from material adjoint contributions into the global objective gradient is consistent in both direction and magnitude in the tested regime.

\section{Case 3 FD vs Adjoint Results}
\subsection{Component Gradients at Initial Guess}
At \(\theta_0=[10,4,0.10]\), adjoint gradients were
\[
g_{\text{adj}}=[-0.31917199,\,-4.48308421,\,10.23837569].
\]
Table~\ref{tab:hsweep} compares FD results across \(h\).

\begin{table}[h!]
\centering
\caption{Case 3 component FD checks at the initial guess. Relative error is
\(\frac{|g_{\mathrm{adj}}-g_{\mathrm{FD}}|}{\max(|g_{\mathrm{adj}}|,|g_{\mathrm{FD}}|,10^{-14})}\).}
\label{tab:hsweep}
\begin{tabular}{@{}lccc@{}}
\toprule
\(h\) & \(g_{\mathrm{FD}}(s_0,s_1,\tau)\) & relative error \((s_0,s_1,\tau)\) \\
\midrule
\(1.0\times10^{-2}\) & \((-0.0125,\ 0.8117,\ 8.1295)\) & \((0.9608,\ 1.1811,\ 0.2060)\) \\
\(3.0\times10^{-3}\) & \((0.0232,\ 0.8144,\ 8.2423)\) & \((1.0727,\ 1.1817,\ 0.1950)\) \\
\(1.0\times10^{-3}\) & \((0.3664,\ 0.5313,\ 7.9278)\) & \((1.8712,\ 1.1185,\ 0.2257)\) \\
\(3.0\times10^{-4}\) & \((-1.5951,\ 0.9229,\ 10.2155)\) & \((0.7999,\ 1.2059,\ 0.0022)\) \\
\(1.0\times10^{-4}\) & \((-1.0430,\ 0.6650,\ 4.8006)\) & \((0.6940,\ 1.1483,\ 0.5311)\) \\
\bottomrule
\end{tabular}
\end{table}

Observation: \(\tau\) can agree at some \(h\) (notably \(h=3\times10^{-4}\)), but both layer gradients remain inconsistent in sign and/or magnitude across all tested \(h\).

\subsection{Directional Derivative Check}
Using direction \(d=[0.84270097,-0.48154341,0.24077171]^T\), the adjoint directional derivative was
\[
g_{\mathrm{adj}}\cdot d = 4.35494430.
\]
FD directional derivatives remained unstable with \(h\):
\begin{table}[h!]
\centering
\caption{Case 3 directional derivative check at \(\theta_0\).}
\label{tab:dir}
\begin{tabular}{@{}ccc@{}}
\toprule
\(h\) & FD directional derivative & relative error \\
\midrule
\(1.0\times10^{-2}\) & \(1.6087\) & \(0.6306\) \\
\(3.0\times10^{-3}\) & \(1.5967\) & \(0.6334\) \\
\(1.0\times10^{-3}\) & \(1.1220\) & \(0.7424\) \\
\(3.0\times10^{-4}\) & \(1.2918\) & \(0.7034\) \\
\(1.0\times10^{-4}\) & \(-3.4515\) & \(1.7925\) \\
\(3.0\times10^{-5}\) & \(12.5019\) & \(0.6517\) \\
\(1.0\times10^{-5}\) & \(8.1136\) & \(0.4633\) \\
\bottomrule
\end{tabular}
\end{table}

Directional FD instability indicates a strongly non-smooth or numerically noisy objective in this regime.

\subsection{Windowed Directional-Derivative Localization (\(h=10^{-2}\))}
To localize where directional mismatch is injected over time, we repeated the directional-derivative comparison on short objective windows rather than only on the full horizon.
Using the same perturbation size \(h=10^{-2}\) and direction \(d=[0.84270097,-0.48154341,0.24077171]^T\), we defined three short windows (early, mid, late), each spanning 200 steps, and one full-horizon reference window.
For each window \(W\), we compared
\[
g_W\cdot d
\quad\text{against}\quad
\frac{J_W(\theta+h d)-J_W(\theta-h d)}{2h},
\]
where \(J_W\) is the window-restricted misfit.

The resulting relative errors were:
\begin{center}
\begin{tabular}{@{}lcccc@{}}
\toprule
Window & time interval & adjoint \(g_W\cdot d\) & FD \(g_W\cdot d\) & relative error \\
\midrule
early & \(t\in[0.00,0.80]\) s & \(1.1815\) & \(0.7913\) & \(0.330\) \\
mid & \(t\in[1.60,2.40]\) s & \(1.5800\) & \(0.3298\) & \(0.791\) \\
late & \(t\in[3.20,4.00]\) s & \(-0.0471\) & \(0.0055\) & \(1.116\) \\
full & \(t\in[0.00,4.00]\) s & \(4.3549\) & \(1.6087\) & \(0.631\) \\
\bottomrule
\end{tabular}
\end{center}

This windowed view indicates that mismatch is not uniformly distributed in time: agreement is comparatively better in the early window and degrades substantially in the mid/late windows, with a sign flip in the late window.
Hence, time-bin localization provides a sharper diagnostic than a single full-horizon number by identifying when directional inconsistency is introduced.

The larger late-horizon mismatch is consistent with dissipative attenuation of wave energy: by late time, the residual signal and corresponding window loss are much smaller than in early time (\(5.92\times10^{-3}\) versus \(1.88\)).
In this low-signal regime, finite-difference quotients are more sensitive to cancellation and resolution limits, while fixed-point rematerialization errors (including repeated overflow-store/load roundoff) represent a larger fraction of the remaining sensitivity.
Accordingly, directional agreement degrades most strongly near the tail of the horizon, where the inverse signal-to-numerical-noise ratio is weakest.

\section{Reduced-Stiffness/Reduced-Horizon Fixed-Point Validation}
To test whether adjoint sensitivities are reliable in a smoother fixed-point regime, we evaluated a control case with reduced stiffness contrast, larger relaxation time, and shorter physical horizon than the high-contrast reference problem.
The validation was performed entirely in fixed-point arithmetic, preserving the same rematerialization-driven computational paradigm.

The control configuration used two layers with \(s_{\text{true}}=[6.0,3.5]\), \(s_{\text{init}}=[4.8,4.2]\), \(\tau_{\text{true}}=1.0\), \(\tau_{\text{init}}=1.3\), and a shortened window of \(N_t=300\) at \(\Delta t=2.5\times10^{-3}\).
At the initial state, the fixed-point adjoint gradient was
\[
g^{\mathrm{fix}}=[-0.02667,\ 4.43059,\ 0.10830].
\]
Using a single, resolution-aware perturbation level \(h=10^{-2}\), the central-difference gradient was
\[
g^{\mathrm{FD}}=[-0.03125,\ 4.46842,\ 0.09738].
\]
The corresponding component-wise relative errors were
\[
[0.1465,\ 0.0085,\ 0.1008],
\]
with worst-component mismatch \(14.65\%\).

As an additional directional check, for a normalized probe direction \(d\), we obtained
\[
g^{\mathrm{fix}}\!\cdot d=-2.1299,\qquad
\left(\frac{J(\theta+h d)-J(\theta-h d)}{2h}\right)_{h=10^{-2}}=-1.9734,
\]
with directional relative error \(7.35\%\).

These results indicate that, in a smoother fixed-point setting, the adjoint delivers gradient magnitudes and directions that are consistent with finite-difference trends.
This supports the validity of the fixed-point adjoint mechanism itself, while remaining consistent with the stronger non-smoothness observed in the full Case 3 high-contrast regime.

\paragraph{Step-size interpretation for finite-difference checks.}
The finite-difference verification error in this dissipative inverse problem is expected to be non-monotone as a function of perturbation size.
For excessively large perturbations, the secant approximation departs from the local linear regime and accumulates truncation bias from objective curvature.
For excessively small perturbations, the objective increment becomes weak relative to effective numerical resolution and damping-induced attenuation, so cancellation and roundoff effects dominate.
Consequently, finite-difference quality is governed by a truncation-versus-resolution tradeoff rather than by monotonic improvement with smaller or larger perturbations.
In the present control regime, the chosen intermediate perturbation level provides the most reliable agreement with adjoint-based directional sensitivity.

\section{Cross-Arithmetic Validation of Gradient Direction}
Because the primary production workflow is fixed-point rematerialized integration, we performed an auxiliary validation in a numerically milder control regime (reduced stiffness contrast, larger \(\tau\), and shorter horizon) to isolate directional consistency of adjoint sensitivities across arithmetic representations.

Let \(g^{\mathrm{fix}}=\nabla_\theta J\) and \(g^{\mathrm{flt}}=\nabla_\theta J\) denote adjoint gradients computed with fixed-point and floating-point viscoelastic updates, respectively, at the same parameter state and under the same objective.
Directional agreement was quantified using:
\[
\mathcal{C}=\frac{g^{\mathrm{fix}}\cdot g^{\mathrm{flt}}}
{\|g^{\mathrm{fix}}\|_2\,\|g^{\mathrm{flt}}\|_2},
\]
where \(\mathcal{C}=1\) indicates perfect collinearity, and by component-wise sign concordance.
As a complementary check, we compared scalar directional derivatives \(g\cdot d\) along a normalized probing direction \(d\).

For the control case, the gradients were
\[
g^{\mathrm{fix}}=[-0.03865,\ 4.82004,\ 2.94000],\qquad
g^{\mathrm{flt}}=[-0.03739,\ 4.89786,\ 2.89158].
\]
The cosine similarity was
\[
\mathcal{C}=0.999896,
\]
with \(100\%\) sign agreement across all gradient components.
For the probing direction \(d\), the directional derivatives were
\[
g^{\mathrm{fix}}\!\cdot d=-1.6458,\qquad
g^{\mathrm{flt}}\!\cdot d=-1.6938,
\]
which corresponds to a cross-arithmetic directional mismatch of \(2.84\%\).

These results show that, in a regime where both arithmetic paths remain well behaved, fixed-point and floating-point adjoint gradients carry essentially the same descent direction.
This supports directional correctness of the fixed-point adjoint signal, while remaining compatible with the observed finite-difference fragility in the full high-contrast Case 3 regime.

\section{Diagnosis and Interpretation}
The evidence supports the following:
\begin{itemize}
\item Core adjoint machinery is likely correct in simplified/unit-tested settings.
\item A dedicated material-point transpose test closes at machine precision (max relative mismatch \(2.38\times10^{-12}\)), supporting constitutive reverse-mode correctness.
\item A dedicated element-level coupling audit on a small fixed-point mesh shows strong adjoint-vs-FD consistency (\(\cos=0.9999993\), relative \(L^2\) mismatch \(7.93\times10^{-3}\), and \(100\%\) sign agreement on active elements).
\item A windowed directional check (\(h=10^{-2}\)) localizes mismatch growth to mid/late portions of the Case 3 horizon (relative errors \(0.791\) and \(1.116\), versus \(0.330\) in early time).
\item A reduced-stiffness/reduced-horizon fixed-point control case with \(\tau=1\) shows materially improved FD consistency (max component relative error \(14.65\%\), directional error \(7.35\%\)).
\item Reduced-control fixed-vs-float cross-check shows strong gradient-direction agreement (\(\cos\approx 0.9999\)).
\item In Case 3 high contrast, \texttt{fixed\_visco} with very frequent overflow checkpointing (\texttt{overflow\_limit=10}, \texttt{mat\_overflow\_limit=10}) produces strong step-size sensitivity.
\item Layer-gradient FD mismatch is persistent (including sign disagreement for \(s_1\)), so Case 3 does \emph{not} currently provide a clean validation of those adjoint components.
\end{itemize}

Hence, the strict answer to ``is adjoint implemented correctly?'' is:
\begin{quote}
\textbf{Partially verified.} The adjoint implementation is verified by dedicated unit tests, but \textbf{not fully verified in Case 3 high-contrast full-scale runs} for layer parameters under the present fixed-point/overflow numerical configuration.
\end{quote}

\section{Recommended Reporting Language (Journal Style)}
Suggested statement for a manuscript:
\begin{quote}
``Adjoint gradients were first validated on unit-test problems at material and system scales, where finite-difference agreement was observed within prescribed tolerances.
For the full high-contrast dissipative-wave benchmark (Case 3), finite-difference verification revealed strong step-size sensitivity and persistent mismatch in layer-wise stiffness gradients, while relaxation-time gradients showed intermittent agreement.
These results indicate that, under the current fixed-point overflow-rematerialization settings, Case 3 does not provide a fully smooth reference for strict finite-difference adjoint validation of all parameters.''
\end{quote}

\section{Reproducibility}
\subsection*{Unit tests}
\begin{verbatim}
cd build
ctest --output-on-failure -R "test_SystemAdjointKick|test_ViscoElasticityAdjoint|test_UniaxialViscoelasticity"
\end{verbatim}

\subsection*{Material-point transpose test (isolated constitutive pullback)}
\begin{verbatim}
cd build
./bin/test_ViscoElasticityAdjoint \
  --gtest_filter=test_ViscoElasticityAdjoint.material_point_transpose_identity_qhistory_stress_seeds
\end{verbatim}

\subsection*{Element-level coupling audit (small-mesh fixed-point control)}
\begin{verbatim}
python examples/dissipative_wave/dissipative_wave_element_coupling_audit.py
\end{verbatim}

\subsection*{Windowed directional-derivative localization (Case 3)}
\begin{verbatim}
python examples/dissipative_wave/dissipative_wave_windowed_directional_check.py
\end{verbatim}

\subsection*{Case 3 reference run}
\begin{verbatim}
python examples/dissipative_wave/dissipative_wave_inverse_cases.py \
  --case case3_high_contrast
\end{verbatim}

\subsection*{Case 3 adjoint diagnostics (script option)}
\begin{verbatim}
python examples/dissipative_wave/dissipative_wave_inverse.py \
  --width 15 --sensor-distribution-width 13.5 \
  --nx 150 --ny 30 --nsteps 1000 --dt 4e-3 \
  --n-layers 2 --n-sensors 5 \
  --true-layers 15,7.7 --init-layers 10,4 \
  --true-tau 0.08 --init-tau 0.10 \
  --integrator-type fixed_visco \
  --overflow-limit 10 --mat-overflow-limit 10 \
  --adjoint-diagnostics --adjoint-diagnostics-fd-check
\end{verbatim}

\subsection*{Reduced-control fixed-vs-float auxiliary cross-check}
\begin{verbatim}
python examples/dissipative_wave/dissipative_wave_aux_control_crosscheck.py
\end{verbatim}

\subsection*{Reduced-control fixed-point validation}
\begin{verbatim}
python examples/dissipative_wave/dissipative_wave_fixed_control_verification.py
\end{verbatim}

\section{Available Figures}
Existing generated figures relevant to this report:
\begin{itemize}
\item \texttt{../case3\_high\_contrast\_summary.png}
\item \texttt{initial\_adjoint\_trace.png}
\item \texttt{recovered\_adjoint\_trace.png}
\item \texttt{truth\_adjoint\_trace.png}
\item \texttt{adjoint\_layer\_grad\_lambda\_compare.png}
\end{itemize}

\end{document}
